\subsection{Three cross states and small opposite corners}

\begin{lemma}[Exact coinvariants with at most three states]
\label{C:three-state-exact-coinvariants}
Let $N$ have projective trivial right rational module, and let
$S$ be an $(N,M)$-biset with $1\leq|S|\leq3$.
If its collage has type $\mathbb{CP}^2$, an ordinary right
$M$-act cone on three states has that type.
In particular $|M|\geq18$.
\end{lemma}
\begin{proof}
The actual orbit-biset functor replaces $N$ by the trivial
monoid and $S$ by the right $M$-set $N\backslash S$.
Exact rational coinvariants give a rational homology and
cohomology-ring equivalence on the edge bars and vertex spaces,
hence on the collages.  If the $N$-action is nontrivial, its
orbit set has at most two states; the finite-line lemma makes
every positive rational product zero.

If the $N$-action is trivial, this collapse has an actual
section: include the trivial monoid into $N$ by its identity
and keep the $M$ vertex and cross arrows unchanged.  Thus the
ordinary act cone is a rationally equivalent homotopy retract
of a hypothetical projective plane.  Its integral homology is
a direct summand in each degree of that of the projective plane.
The rational equivalence forces each such split map to be an
integral isomorphism, since those groups are either zero or
$\mathbb Z$.  Its fundamental group is also trivial, being a
retract of the trivial group.  Homological Whitehead gives an
actual homotopy equivalence with the projective plane.
One or two states are excluded by the finite-line lemma;
Theorem~\ref{C:three-state-cone-minimum22} gives $|M|\geq18$
in the remaining case.
\end{proof}

\begin{proposition}[The five-element sphere corner on three states]
\label{C:P5-three-state-bisets}
Let $P=\{1\}\sqcup(\{0,1\}\times\{0,1\})$ be the
five-element sphere monoid, let $N$ be any finite monoid, and
let $S$ be a nonempty $(P,N)$-biset with $|S|\leq3$.
If $H^2(BN;\mathbb Q)=0$, its collage cannot have type
$\mathbb{CP}^2$.  When $|S|=3$ and the left $P$-action has
no fixed state, the same conclusion holds without the
hypothesis on $H^2(BN;\mathbb Q)$.
\end{proposition}
\begin{proof}
One state gives a join; two states are covered by
Lemma~\ref{C:two-state-biset-H2zero}.  Suppose there are three
states.  If one is fixed by $P$, it gives an actual section
$BP\to E_S$ of the edge projection to $BP$.  If
$H^2(BN;\mathbb Q)=0$, the complete vertex degree-two group
therefore injects into the edge group.  All degree-two classes
of the collage are relative to its vertices, and their products
vanish because that relative quotient is a suspension.

It remains to prove the stronger assertion with no $P$-fixed
state.  Write the band elements as $k_{ij}$.  Each left orbit
of the band in its image has one or two states; a one-point
orbit would be a fixed state.  There is therefore one two-point
core $U=\{u_0,u_1\}$ and a third state $z$.  The idempotent
identity $k_{ij}(k_{ij}s)=k_{ij}s$ forces each band actor to
be constant on all three states, with
\[
 k_{ij}s=u_i\qquad(s\in S).
\]
Commutation makes every right $N$-actor fix both $u_0,u_1$.
It can only fix $z$ or send it to one of these core states.
Define $\chi_n=1$ in the first case and $\chi_n=0$ otherwise;
put $\beta_n=0$ when $\chi_n=1$, and put $\beta_n=i$ when
$zn=u_i$.  The action equations give
\begin{equation}\label{C:P5-three-state-affine}
 \chi_{mn}=\chi_m\chi_n,\qquad
 \beta_{mn}=\beta_m+\chi_m\beta_n.
\end{equation}

We need the actual equivariant model of the left action bar
$Z=B(P\ltimes S)$.  The subcategory on $U$ has contractible
classifying space.  Indeed the arrows $k_{i0}:u_0\to u_i$
give a natural transformation from the constant functor at
$u_0$ to its identity functor.  Every arrow leaving $z$ has
label $k_{ij}$, and postcomposition preserves $j$.  Thus its
nerve is the union of two cones on the core action category,
one for each column $j$, with their cone vertices both equal
to $z$.  Their intersection consists of the core category
and the singleton $z$.

Collapse the contractible core category to one vertex.  This
gives an equivalence to the graph with two vertices and two
parallel edges, hence a circle.  It is $N$-equivariant: if
$zn=z$, the actor is the identity on the whole action category;
otherwise it sends the whole category into its core, so it
acts constantly on the quotient graph.  In particular its
action on the circle's first homology is $\chi$.
The quotient circle has a fixed core vertex, and its cellular
chain complex splits over $\mathbb QN$ as
\[
 \mathbb Q\ \oplus\ \mathbb Q_\chi[1]
\]
plus a contractible two-term summand.  The edge double bar
$E_S=Z_{hN}$ therefore satisfies, compatibly with projection
and its core-state section,
\begin{equation}\label{C:P5-three-state-edge-homology}
 H_n(E_S;\mathbb Q)\cong H_n(BN;\mathbb Q)
       \oplus\operatorname{Tor}^{\mathbb QN}_{n-1}
                     (\mathbb Q_\chi,\mathbb Q)
 \qquad(n\geq1).
\end{equation}

If $\chi$ is identically one, the $N$-action is trivial and
$E_S=BN\times Z$.  Its map to $BP\simeq S^2$ factors through
$Z\simeq S^1$ and is null-homotopic.  Since a finite monoid
has zero first rational homology, Mayer--Vietoris gives
$\dim H_2(BC;\mathbb Q)=2$, excluding a projective plane.

Suppose $\chi$ is nontrivial.  Its coinvariants are zero, so
$H_1(E_S;\mathbb Q)=0$.  The section in
\eqref{C:P5-three-state-edge-homology} maps constantly to $BP$.
After cancelling its $BN$ summands, the relevant
Mayer--Vietoris segment is therefore
\[
 0\longrightarrow H_3(BC;\mathbb Q)
 \longrightarrow T_1\xrightarrow{\alpha}\mathbb Q
 \longrightarrow H_2(BC;\mathbb Q)\longrightarrow0,
 \qquad
 T_1=\operatorname{Tor}^{\mathbb QN}_1
                 (\mathbb Q_\chi,\mathbb Q).
\]
A projective plane would force $\alpha=0$ to retain its
one-dimensional second homology, and then force $T_1=0$
by its vanishing third homology.

Let $K$ be the minimal ideal of $N$, with row set $I$.
The character $\chi$ is zero on $K$: the ideal generated by
any zero-character element contains $K$.  Averaging the finite
minimal maximal group gives a projective left row module
$P_I=\mathbb QI$ and an exact sequence
\[
 0\longrightarrow V\longrightarrow P_I
   \longrightarrow\mathbb Q\longrightarrow0,
 \qquad V=\widetilde{\mathbb QI}.
\]
Since $P_I=KP_I$, tensoring with $\mathbb Q_\chi$ kills
$P_I$.  Projectivity and dimension shifting give
\[
 T_1\cong\mathbb Q_\chi\otimes_{\mathbb QN}V.
\]
Equation~\eqref{C:P5-three-state-affine} makes $\beta$ constant
on each minimal right ideal, so write its value on row $i$
as $b_i$.  Left multiplication by $n$ induces the row map
$\sigma_n$, and the same equation gives
\[
 b_{\sigma_n(i)}=\beta_n+\chi_n b_i.
\]
If the $b_i$ are not all equal, their differences define a
nonzero module homomorphism
$V\to\mathbb Q_\chi$.  Tensoring it gives a surjection
$T_1\to\mathbb Q$, contradicting $T_1=0$.

Consequently $b_i=c$ for a single $c\in\{0,1\}$, and the
last equation gives $\beta_n=c(1-\chi_n)$ for every $n$.
Every nonidentity action on $z$ therefore collapses it to the
same core state $u_c$.  The map $S\to U$ fixing $U$ and
sending $z$ to $u_c$ is an actual biset retraction.
The collage on $U$ is equivalent to $BP$: its edge is the
product of the contractible core action bar with $BN$, and
its projection to $BN$ is an equivalence.
The original collage thus retracts onto $S^2$.  If its second
rational cohomology had dimension one, the pulled-back sphere
class would span it and have square zero.  This again excludes
the projective plane, completing every case.
\end{proof}

\begin{corollary}\label{C:sixteen-P5-opposite-eight}
In the sixteen-arrow triangular two-object calculation, the
profile $(5,8,3)$ with rationally degree-two acyclic
eight-element corner is excluded structurally.  If that corner
has nonzero rational $H^2$, opposition places the case in
the complete $(8,5,3)$ catalog.
No exhaustive catalog of degree-two acyclic eight-element
monoids is required for this reduction.
\end{corollary}
\begin{proof}
The only five-element monoid with nonzero rational $H^2$
is the sphere monoid $P$, which is isomorphic to its opposite.
Apply the preceding proposition and then take the opposite
category when appropriate.
\end{proof}
